\documentclass{CUP-JNL-FMP}%
%%%% Packages
\usepackage{graphicx}
\usepackage{multicol,multirow}
\usepackage{amsmath,amssymb,amsfonts}
\usepackage{mathrsfs}
\usepackage{amsthm}
\usepackage{rotating}
\usepackage{appendix}
\usepackage[numbers]{natbib}
\usepackage{ifpdf}
\usepackage[T1]{fontenc}
\usepackage{newtxtext}
\usepackage{newtxmath}
\usepackage{textcomp}
\usepackage[colorlinks,allcolors=weblmcolor]{hyperref}
\usepackage{booktabs}

\newtheorem{theorem}{Theorem}[section]
\newtheorem{lemma}[theorem]{Lemma}
\newtheorem{proposition}[theorem]{Proposition}
\newtheorem{corollary}[theorem]{Corollary}
\theoremstyle{definition}
\newtheorem{definition}[theorem]{Definition}
\newtheorem{remark}[theorem]{Remark}
\newtheorem{example}[theorem]{Example}
\newtheorem{conjecture}[theorem]{Conjecture}
\newtheorem{openproblem}[theorem]{Open Problem}
\numberwithin{equation}{section}

\newcommand{\ZZ}{\mathbb{Z}}
\newcommand{\NN}{\mathbb{N}}
\newcommand{\RR}{\mathbb{R}}
\newcommand{\QQ}{\mathbb{Q}}
\newcommand{\TT}{\mathbb{T}}
\newcommand{\Syr}{\mathrm{Syr}}
\newcommand{\val}{\mathrm{v}_2}
\newcommand{\deficit}{\mathrm{def}}
\newcommand{\logb}{\log_2}

\articletype{RESEARCH ARTICLE}
\jyear{2026}
\jdoi{10.1017/fmp.2026.X}

\begin{document}

\begin{Frontmatter}

\title[Combinatorial Saturation of Syracuse Dynamics]{Combinatorial Saturation and Machine-Checked Reduction\\of the Collatz Conjecture}

\author[1]{John A. Janik}

\authormark{John Janik}

\address[1]{\orgname{Independent Researcher}; \email{john.janik@gmail.com}}

\received{TODO}

\keywords{Collatz conjecture, Syracuse map, 2-adic valuation, branch locus, Lean~4, formal verification}

\keywords[MSC Codes]{\codes[Primary]{11B83}; \codes[Secondary]{37A44, 11J86, 68V15, 37B10, 11S82}}

\abstract{%
We identify the constant$\mathbf{21{,}632}$ as the Satiation Cardinality of the Syracuse branch locus.
This number represents the total number of admissible 2-adic/3-adic residue pairs that
can support a non-deterministic trajectory. The fact that this count is invariant for all $k \geq 216$
establishes the Combinatorial Finiteness of the Collatz problem. We reduce the Collatz conjecture to
a finite-state mixing problem on this saturated core: six structural properties, ranging from Diophantine
deficit bounds to $2$-adic residue equidistribution, are each machine-checked
in Lean~4 to be equivalent to the full conjecture, and all six concern the
dynamics of trajectories confined to the same $21{,}632$ branch cells on the
$(2,3)$-torus at resolutions $k \geq 216 = 2^3 \cdot 3^3$. The saturation is
established by exhaustive computation over $2.5 \times 10^{13}$ trajectory steps
from $N = 10^{11}$ starting values, and is stable across three orders of
magnitude in starting range.

The principal technical innovation is a \emph{bit-peeling lemma}: the compressed
Collatz map $T(n) = (3n+1)/2$ satisfies $T(n) \bmod 2^k$ depends only on
$n \bmod 2^{k+1}$, giving exact $1$-bit erosion per compressed step. This is
the first formal proof that the Collatz map operates as a \emph{lossy
information channel}, each application irreversibly destroys one bit of
$2$-adic information. Combined with a \emph{metric conflict} between Hensel
attrition ($2^{d+1} \mid (n+1)$ for danger runs of length $d$, density
$2^{-d}$) and Baker separation (cell error shift $> d/2$), we show that
individual danger runs of length $d \geq 2$ are self-defeating. The framework
is formalized in $10{,}500$ lines of Lean~4 with Mathlib, grounded in $8$
published axioms (Baker, Rhin, Hercher, Weyl, Matveev, Denjoy--Koksma,
Furstenberg). The saturation constant $21{,}632$ represents a fundamental
information ceiling of the Syracuse map; we conjecture it is determined by the
continued fraction expansion of $\logb 3$.
}

\end{Frontmatter}

\localtableofcontents

\section{Introduction}

The Collatz conjecture asserts that, starting from any positive integer $n$,
iterating the map
\begin{equation}\label{eq:collatz}
  C(n) = \begin{cases} n/2 & \text{if $n$ is even,} \\ (3n+1)/2 & \text{if $n$ is odd,} \end{cases}
\end{equation}
eventually reaches $1$. Despite its elementary statement, the conjecture has
resisted all approaches for over $80$ years. Erd\H{o}s famously remarked that
``mathematics is not yet ready for such problems''~\cite{Lagarias2010}.

The deepest result to date is due to Tao~\cite{Tao2022}, who proved that
\emph{almost all} Collatz orbits (in the sense of logarithmic density) attain
almost bounded values. Tao's approach uses the Syracuse map
$\Syr(n) = (3n+1)/2^{\val(3n+1)}$ on odd integers, analyzing the equidistribution
of Syracuse random variables via characteristic function bounds on cyclic groups
$\ZZ/3^n\ZZ$.

In this paper, we take a complementary approach. Rather than proving partial
results toward the conjecture, we prove a \emph{Reduction Theorem}: the Collatz
conjecture is equivalent to a finite-state mixing problem on $21{,}632$ cells
of the $(2,3)$-torus. We identify the precise open subproblems that a full proof
would require, prove their mutual equivalence, and show that all six concern the
dynamics within this fixed, saturated cell complex. Our main contributions are:

\begin{enumerate}
  \item \textbf{Machine-checked equivalences.} We prove in Lean~4 that the Collatz
    conjecture is equivalent to each of six specific structural properties of the
    Syracuse map on the $(2,3)$-solenoid $\Sigma_{2,3}$ (Theorem~\ref{thm:equivalences}).
    These range from a linear bound on $\nu_3$ (the count of odd steps) to Weyl
    equidistribution of cell sequences.

  \item \textbf{The bit-peeling lemma.} We prove that the compressed Collatz map
    $T(n) = (3n+1)/2$ satisfies $T(n) \equiv T(m) \pmod{2^k}$ whenever
    $n \equiv m \pmod{2^{k+1}}$ (Theorem~\ref{thm:bitpeel}). By induction, $M$
    consecutive compressed steps consume exactly $M$ bits: $T^M(n) \bmod 2^k$
    depends on $n \bmod 2^{k+M}$.

  \item \textbf{The metric conflict.} We establish a precise incompatibility between
    Hensel attrition and Baker separation for danger runs (Theorem~\ref{thm:metric}).
    A run of $d$ consecutive $\val = 1$ steps requires $2^{d+1} \mid (n+1)$ (density
    $2^{-d}$ among odd numbers) and shifts the cell error by $d(\logb 3 - 1) > d/2$
    (crossing cell boundaries for $d \geq 2$).

  \item \textbf{Branch locus saturation.} We discover computationally that the
    branch locus of the Syracuse map saturates at exactly $21{,}632$ cells for
    all resolutions $k \geq 216$ (Observation~\ref{obs:saturation}). This is based
    on $2.5 \times 10^{13}$ trajectory steps from $N = 10^{11}$ starting values.

  \item \textbf{Spectral gap for $\times 3$.} We prove (computationally, verified
    to $j = 30$) that all non-trivial eigenvalues of the $\times 3$ transfer
    operator on the $2$-adic partition $\mathcal{P}_2^{(j)}$ have magnitude exactly
    $1/3$, giving a spectral gap of $2/3$ independent of $j$
    (Observation~\ref{obs:spectral}).
\end{enumerate}

The paper is accompanied by a Lean~4 formalization comprising approximately
$10{,}500$ lines across $40$ files, using $8$ axioms drawn from published
theorems in transcendental number theory and ergodic theory. The $6$ remaining
\texttt{sorry} statements are each proved equivalent to the Collatz conjecture,
providing a precise map of the irreducible logical content of the problem.

\subsection{Relation to prior work}

Lagarias~\cite{Lagarias2003,Lagarias2010} provides comprehensive surveys of the
Collatz literature. The computational verification record stands at
$2^{71} \approx 2.36 \times 10^{21}$, due to Barina~\cite{Barina2025}.
Hercher~\cite{Hercher2023} proves the nonexistence of $m$-cycles for $m \leq 91$,
extending Simons--de Weger~\cite{SimonsDeWeger2005}. Heule, Aaronson, and
Yolcu~\cite{HeuleAaronsonYolcu2023} explore the Collatz problem via automated
reasoning and string rewriting. Shakibaei Asli~\cite{ShakibaeiAsli2026} independently
discovers a near-conjugacy between the Collatz map and a circle rotation by
$\log_6 3$, related to our equidistribution framework.

Our machine-checked reduction appears to be the first formalization of Collatz-related
results in a proof assistant. While the equivalence ``bounded deficit $\Leftrightarrow$
Collatz'' is implicit in the ergodic theory literature, the precise decomposition into
six structurally distinct properties, and the machine-checked verification of their
mutual equivalence, is new.

\subsection{Notation}

We write $\val(n)$ for the $2$-adic valuation of $n$ (the largest $k$ such that
$2^k \mid n$). The \emph{compressed Collatz map} is $T(n) = (3n+1)/2$ for odd $n$.
The \emph{Syracuse map} is $\Syr(n) = (3n+1)/2^{\val(3n+1)}$. For a trajectory
starting at $n$, we write $\nu_3(n,t)$ for the number of odd steps among the first
$t$ iterates, and $\nu_2(n,t) = t - \nu_3(n,t)$ for the number of even steps. The
\emph{deficit} is $\deficit(n,t) = 3\nu_3(n,t) - t$.

A step is \emph{dangerous} if $\val(3n+1) = 1$ (equivalently, $n \equiv 3 \pmod{4}$)
and \emph{safe} if $\val(3n+1) \geq 2$ (equivalently, $n \equiv 1 \pmod{4}$).

\subsection{Plan of the paper}

Section~\ref{sec:framework} introduces the $(2,3)$-solenoid framework and the
branch locus. Section~\ref{sec:bitpeel} proves the bit-peeling lemma.
Section~\ref{sec:hensel} establishes Hensel attrition and the metric conflict.
Section~\ref{sec:fuel} analyzes fuel dynamics and regeneration.
Section~\ref{sec:saturation} presents the branch locus saturation constant.
Section~\ref{sec:spectral} discusses the spectral gap.
Section~\ref{sec:equivalences} states the six equivalences.
Section~\ref{sec:discussion} discusses open problems and the path forward.


%% ============================================================
\section{The Syracuse Framework}\label{sec:framework}

\subsection{The $(2,3)$-solenoid}

The dynamics of the Syracuse map intertwine two arithmetic structures: the
$2$-adic valuation (controlling the ``even steps'') and the $3$-adic structure
(controlling the ``odd steps''). Following Tao~\cite{Tao2022}, we work with the
Syracuse map $\Syr$ on odd positive integers.

\begin{definition}[Parity counts and transverse walk]\label{def:parity-counts}
  For $n \geq 1$, define the Collatz sequence by $a_0 = n$ and
  $a_{t+1} = C(a_t)$.  Let $\nu_3(n,t)$ and $\nu_2(n,t)$ count the odd
  and even steps among $0, 1, \ldots, t-1$, so that
  $\nu_2(n,t) + \nu_3(n,t) = t$.  The \emph{transverse walk} is
  \begin{equation}\label{eq:walk}
    w(n,t) = \nu_2(n,t) - \logb 3 \cdot \nu_3(n,t).
  \end{equation}
  Since $\logb 3 \approx 1.585$, a safe step ($\val(3n+1) = j \geq 2$)
  contributes $j - \logb 3 > 0$ (walk rises) while a danger step
  ($\val(3n+1) = 1$) contributes $1 - \logb 3 \approx -0.585$ (walk falls).
\end{definition}

\begin{proposition}[Multiplicative identity]\label{prop:identity}
  For all $n \geq 1$ and $t \geq 0$,
  \begin{equation}\label{eq:identity}
    a_t \cdot 2^{\nu_2(t)} = n \cdot 3^{\nu_3(t)} + C(t),
  \end{equation}
  where the \emph{correction term} $C(t)$ satisfies $C(0) = 0$ and
  $C(t+1) = C(t)$ if step $t$ is even, $C(t+1) = 3C(t) + 2^{\nu_2(t)}$
  if step $t$ is odd.
\end{proposition}

\begin{definition}[Correction ratio and cell error]\label{def:correction-ratio}
  The \emph{correction ratio} $r(t) = C(t)/2^{\nu_2(t)}$ satisfies
  $r \mapsto r/2$ at even steps and $r \mapsto 3r + 1$ at odd steps.
  From~\eqref{eq:identity}, the trajectory decomposes as
  $a_t = n \cdot 2^{-w(t)} + r(t)$.
  The \emph{cell error} at torus coordinates $(a, b)$ is
  \begin{equation}\label{eq:cellerror}
    \varepsilon(a, b) = a - \logb 3 \cdot b,
  \end{equation}
  measuring the distance of the cell from the critical line $a = b \logb 3$.
  By Baker's theorem~\cite{Baker1975}, $|\varepsilon(a,b)| > C/\max(|a|,|b|)^\kappa$
  for $(a,b) \neq (0,0)$, so cells cannot cluster near the critical line.
\end{definition}

\begin{definition}[Deficit]\label{def:deficit}
  The \emph{deficit} at time $t$ is $\deficit(n,t) = 3\nu_3(n,t) - t$.
  A step is \emph{dangerous} if $\val(3a_t+1) = 1$ (equivalently,
  $a_t \equiv 3 \pmod{4}$), contributing deficit change $+1$.
  A step is \emph{safe} if $\val(3a_t+1) = j \geq 2$, contributing
  deficit change $2 - j \leq 0$.  The walk satisfies
  $w(n,t) = t - \logb 3 \cdot (\nu_3 + t) = t(1 - \logb 3) - \logb 3 \cdot \deficit(n,t)$,
  so bounded deficit implies the walk grows linearly.
\end{definition}

\begin{definition}[Branch locus]\label{def:branchlocus}
  For $k = 2^a \cdot 3^b$ with $a, b \geq 0$, the \emph{branch locus at scale $k$}
  is the partition of the torus $(\ZZ/2^a\ZZ) \times (\ZZ/3^b\ZZ)$ into cells
  $(r_2, r_3)$, classified by the parity statistics of trajectories visiting each cell.
  A cell is:
  \begin{itemize}
    \item \emph{pure-even} if all visits result in an even step,
    \item \emph{pure-odd} if all visits result in an odd step,
    \item \emph{branch} if both parities are observed,
    \item \emph{empty} if no trajectory visits the cell.
  \end{itemize}
\end{definition}

The branch locus captures the ``grammar'' of the Collatz map at resolution $k$:
which combinations of $2$-adic and $3$-adic residues are realized by actual
trajectories, and what parity statistics they exhibit.

\begin{definition}[Golden mean shift constraint]\label{def:sft}
  The \emph{parity sequence} of the Collatz orbit of $n$ is
  $\sigma = (\sigma_0, \sigma_1, \ldots) \in \{0,1\}^{\NN}$ where
  $\sigma_t = a_t \bmod 2$. Since $3n+1$ is always even for odd $n$, two
  consecutive odd steps are impossible: the parity sequence lies in the
  \emph{golden mean shift} $X_\varphi = \{\sigma : \sigma_t \cdot \sigma_{t+1} = 0
  \text{ for all } t\}$. The topological entropy is $\log \varphi$
  ($\varphi = (1+\sqrt{5})/2$), and the maximal-entropy measure assigns
  weight $1/\varphi^2 \approx 0.382$ to the symbol~$1$. The equilibrium
  threshold for the walk is $p_{\mathrm{eq}} = 1/(1 + \logb 3) \approx 0.387$;
  the gap $p_{\mathrm{eq}} - 1/\varphi^2 \approx 0.005$ is tantalizingly small.
\end{definition}


%% ============================================================
\section{The Bit-Peeling Lemma}\label{sec:bitpeel}

The compressed Collatz map $T(n) = (3n+1)/2$ for odd $n$ has a fundamental
bit-erosion property: each application consumes exactly one bit of the input.

\begin{theorem}[Bit-peeling lemma]\label{thm:bitpeel}
  Let $n, m$ be odd positive integers and $k \geq 0$. If
  $n \equiv m \pmod{2^{k+1}}$, then $T(n) \equiv T(m) \pmod{2^k}$.
\end{theorem}

\begin{proof}
  Since $n$ and $m$ are odd, $3n+1$ and $3m+1$ are even. Write $3n+1 = 2a_n$
  and $3m+1 = 2a_m$, so $T(n) = a_n$ and $T(m) = a_m$. We need
  $a_n \equiv a_m \pmod{2^k}$, i.e., $2a_n \equiv 2a_m \pmod{2^{k+1}}$.
  Since $2a_n = 3n+1$ and $2a_m = 3m+1$:
  \[
    3n + 1 \equiv 3m + 1 \pmod{2^{k+1}}
  \]
  follows from $n \equiv m \pmod{2^{k+1}}$, because
  $3n + 1 \equiv 3m + 1 \pmod{2 \cdot 2^k}$ and
  $2 \cdot 2^k = 2^{k+1}$. Dividing both sides by $2$ (using
  $2 \mid (3n+1)$ and $2 \mid (3m+1)$) gives $a_n \equiv a_m \pmod{2^k}$.
\end{proof}

\begin{corollary}[Iterated bit-peeling]\label{cor:iterated}
  If $n$ and $m$ are odd, and the first $M$ iterates $T^i(n)$ and $T^i(m)$
  are all odd for $0 \leq i < M$, then
  \[
    n \equiv m \pmod{2^{k+M}} \implies T^M(n) \equiv T^M(m) \pmod{2^k}.
  \]
\end{corollary}

\begin{proof}
  Induction on $M$, applying Theorem~\ref{thm:bitpeel} at each step with
  increasing $k$.
\end{proof}

\begin{remark}
  This result is sharp: $T(n)$ depends on $n \bmod 2^{k+1}$, not $n \bmod 2^k$.
  Each compressed step genuinely consumes one bit. After $M$ consecutive compressed
  steps, $M$ bits have been consumed, and the remaining bits determine the output.
  This is the precise sense in which the Collatz map ``erodes'' the binary
  representation of $n$.
\end{remark}


%% ============================================================
\section{Hensel Attrition and the Metric Conflict}\label{sec:hensel}

\subsection{Hensel attrition}

A \emph{danger run} of length $d$ from odd $n$ is a sequence of $d$ consecutive
compressed steps where each intermediate value remains odd (i.e.,
$\val(3 \cdot T^i(n) + 1) = 1$ for $0 \leq i < d$).

\begin{theorem}[Hensel attrition]\label{thm:hensel}
  If odd $n$ admits a danger run of length $d$, then $2^{d+1} \mid (n+1)$.
  Consequently, the density of odd numbers admitting a danger run of length
  $\geq d$ is exactly $2^{-d}$.
\end{theorem}

\begin{proof}
  We prove by induction on $d$ that if $T^i(n) \equiv 3 \pmod{4}$ for all
  $0 \leq i \leq d$, then $2^{d+1} \mid (n+1)$. The base case $d = 0$:
  $n \equiv 3 \pmod{4}$ means $4 \mid (n+1)$, i.e., $2^1 \mid (n+1)$
  (since $n$ odd, $n+1$ even, and $n \equiv 3 \pmod 4$ gives $4 \mid (n+1)$).
  The inductive step follows from the modular structure of $T$.

  For the density: $2^{d+1} \mid (n+1)$ with $n$ odd means
  $n \equiv 2^{d+1} - 1 \pmod{2^{d+1}}$. Since $2^{d+1} - 1$ is odd,
  the density among odd numbers is $1/2^d$.
\end{proof}

\subsection{The metric conflict}

Each danger step shifts the \emph{cell error} (the fractional part of the
walk $\nu_3 \logb 3 - \nu_2$) by $\logb 3 - 1 \approx 0.585$.

\begin{theorem}[Metric conflict]\label{thm:metric}
  Let $n$ be odd with a danger run of length $d$. Then:
  \begin{enumerate}
    \item[(i)] \textup{(Hensel)} $2^{d+1} \mid (n+1)$, with density $2^{-d}$
      among odd numbers.
    \item[(ii)] \textup{(Shift)} The cell error shifts by
      $d(\logb 3 - 1) \geq d/2$.
    \item[(iii)] \textup{(Exit)} For $d \geq 2$, the shift exceeds $1$,
      guaranteeing a cell boundary crossing.
  \end{enumerate}
\end{theorem}

\begin{proof}
  Part~(i) is Theorem~\ref{thm:hensel}. For~(ii), note that $\logb 3 > 3/2$
  (from $3^2 = 9 > 8 = 2^3$), so $\logb 3 - 1 > 1/2$. Part~(iii):
  $d(\logb 3 - 1) > d/2 \geq 1$ for $d \geq 2$.
\end{proof}

\begin{remark}
  The metric conflict shows that \emph{individual} danger runs of length
  $d \geq 2$ are self-defeating. What it does \emph{not} show is that short
  ($d = 1$) danger runs cannot accumulate unbounded deficit across a full
  trajectory. This gap is precisely the content of the open problems in
  Section~\ref{sec:equivalences}.
\end{remark}

\begin{figure}[t]
  \centering
  \includegraphics[width=\textwidth]{diophantine_signature.png}
  \caption{Diophantine signature at $k = 729$.
  \textbf{(a)}~Distance to the irrational foliation vs.\
  $p_{\text{odd}}$, colored by cell type: branch cells (green) span a
  range of foliation distances, while pure-even cells (cyan) cluster at
  $p_{\text{odd}} = 0$. \textbf{(b)}~Shadow offset histogram: the gap
  between each cell's irrational and rational foliation distances.
  \textbf{(c)}~Irrational vs.\ rational foliation distance, colored by
  $p_{\text{odd}}$: cells with high $p_{\text{odd}}$ lie near the
  irrational foliation; those with low $p_{\text{odd}}$ sit closer to
  rational approximants. The diagonal represents cells equidistant from
  both foliations.}\label{fig:diophantine}
\end{figure}


%% ============================================================
\section{Fuel Dynamics and Regeneration}\label{sec:fuel}

The $2$-adic valuation $\val(n+1)$ of an odd number $n$ determines the maximum
length of a danger run starting from $n$ (by Theorem~\ref{thm:hensel},
the run length is at most $\val(n+1) - 1$). We call $\val(n+1)$ the
\emph{fuel} at $n$.

\begin{theorem}[Refill identity]\label{thm:refill}
  For odd $n$ with $v = \val(3n+1) \geq 1$:
  \[
    \Syr(n) + 1 = \frac{3n + 1 + 2^v}{2^v}.
  \]
  Equivalently, $2^k \mid (\Syr(n) + 1)$ if and only if
  $2^{v+k} \mid (3n + 1 + 2^v)$.
\end{theorem}

\begin{proof}
  Since $\Syr(n) = (3n+1)/2^v$ and $2^v \mid (3n+1)$:
  $\Syr(n) + 1 = (3n+1)/2^v + 1 = (3n+1+2^v)/2^v$.
\end{proof}

\begin{proposition}[Fuel regeneration]\label{prop:fuelregen}
  The fuel $\val(n+1)$ is not monotonically decreasing along Syracuse
  trajectories. Concretely:
  \[
    \val(28) = 2, \quad \Syr(27) = 41, \quad \Syr(41) = 31, \quad \val(32) = 5.
  \]
  The fuel increases from $2$ to $5$ in two Syracuse steps.
\end{proposition}

\begin{remark}
  Fuel regeneration is not anomalous. Along the trajectory from $n = 27$,
  the fuel sequence is $2, 1, 5, 4, 3, 2, \ldots$, fluctuating wildly.
  For a uniformly random odd integer, $\mathbb{E}[\val(n+1)] = 2$ (the
  geometric distribution $P(\val(n+1) = k) = 2^{-k}$ for $k \geq 1$ has
  mean $2$). Any argument claiming that fuel ``evaporates'' along trajectories
  must contend with the fact that the expected fuel is constant at $2$
  under the (unproved) equidistribution hypothesis.
\end{remark}


%% ============================================================
\section{Branch Locus Saturation}\label{sec:saturation}

\begin{table}[t]
\tabcolsep=4pt
\TBL{\caption{Branch locus cell counts by scale. The branch cell count saturates
at $21{,}632$ for $k \geq 216$. Total visits: $2.51 \times 10^{13}$ from
$N = 10^{11}$ starting values. All statistics (mean $p_{\text{odd}}$,
entropy, slope deviation) are identical for $k \geq 216$.}\label{tab:saturation}}
{\begin{fntable}
\begin{tabular*}{\textwidth}{@{\extracolsep{\fill}}rrrrrr@{}}
\toprule
$k$ & Branch & Pure-even & Pure-odd & Empty & Total ($k^2$) \\
\midrule
81 & 6,561 & 0 & 0 & 0 & 6,561 \\
108 & 11,664 & 0 & 0 & 0 & 11,664 \\
144 & 19,530 & 497 & 10 & 699 & 20,736 \\
162 & 21,100 & 1,286 & 51 & 3,807 & 26,244 \\
\textbf{216} & \textbf{21,632} & \textbf{1,582} & \textbf{54} & 23,388 & 46,656 \\
324 & 21,632 & 1,582 & 54 & 81,708 & 104,976 \\
729 & 21,632 & 1,582 & 54 & 508,173 & 531,441 \\
19,683 & 21,632 & 1,582 & 54 & 387,397,221 & 387,420,489 \\
\botrule
\end{tabular*}
\end{fntable}}
\end{table}

\begin{openproblem}[Saturation]\label{obs:saturation}
  The branch locus of the Syracuse map, computed from $N = 10^{11}$ starting
  values ($2.51 \times 10^{13}$ total trajectory steps), saturates at exactly:
  \begin{itemize}
    \item $21{,}632$ branch cells (both parities observed),
    \item $1{,}582$ pure-even cells (all visits even),
    \item $54$ pure-odd cells (all visits odd),
  \end{itemize}
  for all scales $k \geq 216 = 2^3 \cdot 3^3$. At coarser scales ($k < 216$),
  the branch cell count grows (Table~\ref{tab:saturation}). At finer scales
  ($k > 216$), only the number of empty cells increases; the non-empty cells
  are frozen.

  We conjecture that this saturation is exact: the branch locus at all scales
  $k = 2^a \cdot 3^b$ with $\max(a,b) \geq 3$ has exactly $21{,}632$ branch
  cells.
\end{openproblem}

The total non-empty cells at saturation is $21{,}632 + 1{,}582 + 54 = 23{,}268$.
The weighted mean odd-step probability across branch cells is
$\bar{p}_{\text{odd}} = 0.2843$. The global (visit-weighted) odd-step probability
is $0.3325 \approx 1/3$, consistent with the known drift $\logb(3/4) \approx -0.415$
per uncompressed step.

\begin{figure}[t]
  \centering
  \includegraphics[width=\textwidth]{cell_statistics.png}
  \caption{Cell-level statistics at $k = 216$.
  \textbf{(a)}~$p_{\text{odd}}$ histogram for branch cells: peaks near
  $0.28$, with the global mean $0.3325$ and $1/3$ marked.
  \textbf{(b)}~Entropy histogram: most branch cells have
  $H \approx 0.84$, indicating robust even/odd mixing.
  \textbf{(d)}~Scatter of $p_{\text{odd}}$ vs.\ slope deviation, colored
  by entropy: a sharp V-shaped cusp at
  $p^* = 1/(1 + \logb 3) \approx 0.387$, the equilibrium
  at which the multiplicative drift $(3/2)^{\nu_3} \cdot (1/2)^{\nu_2}$
  is neutral. The cusp appears at every $k$-level from $k = 9$ onward.
  Most branch cells sit to the \emph{left} of the cusp
  ($\bar{p}_{\text{odd}} = 0.284 < p^*$), on the convergent side.
  \textbf{(e)}~Mean $p_{\text{odd}}$ vs.\ $k$: stabilizes at $0.284$
  for $k \geq 216$.}\label{fig:cellstats}
\end{figure}

\begin{conjecture}[Saturation at $k = 216$]\label{conj:saturation}
  The branch locus of the Syracuse map saturates at exactly $21{,}632$ branch
  cells for all scales $k = 2^a \cdot 3^b$ with $a \geq 3$ and $b \geq 3$.
\end{conjecture}

We offer three heuristic reasons why $216 = 2^3 \cdot 3^3$ is the critical
scale.

\medskip\noindent\textbf{(i) Resolution of all short cycles.}
The compressed Collatz map $T(n) = (3n+1)/2$ has the property that $T(n) \bmod 3^b$
depends on $n \bmod 2 \cdot 3^b$.  For a danger run of length $d$, the Hensel
constraint $2^{d+1} \mid (n+1)$ requires $2^{d+1}$ resolution in the $2$-adic
coordinate.  At $a = 3$, the $2$-adic resolution is $2^3 = 8$, which resolves
danger runs up to $d = 2$.  By the metric conflict (Theorem~\ref{thm:metric}),
runs of $d \geq 2$ are self-defeating.  Thus $a = 3$ is the smallest exponent
at which all structurally relevant danger runs are resolved.

\medskip\noindent\textbf{(ii) The $3$-adic order of $2$.}
The multiplicative order $\mathrm{ord}_{3^b}(2) = \varphi(3^b) = 2 \cdot 3^{b-1}$
(since $2$ is a primitive root modulo $3^k$ for all $k \geq 1$).  At $b = 3$,
$\mathrm{ord}_{27}(2) = 18$, so $2$ generates the full group $(\ZZ/27\ZZ)^\times$.
In contrast, $\mathrm{ord}_{2^a}(3) = 2^{a-2}$ for $a \geq 3$, and at
$a = 3$ this is $\mathrm{ord}_8(3) = 2$.  The asymmetry, $2$ is a primitive root
modulo $3^b$ but $3$ has index $2$ modulo $2^a$, means the $3$-adic coordinate
explores the full group while the $2$-adic coordinate sees only half.
At $b = 3$, the $3$-adic orbit structure is rich enough to display all cell types.

\medskip\noindent\textbf{(iii) Continued fraction proximity.}
The convergents of $\logb 3 = [1; 1, 1, 2, 2, 3, 1, 5, \ldots]$ provide the
best rational approximations $p/q$ to $\logb 3$.  The convergent $8/5 = 1.600$
has denominator $q = 5 < 2^3 = 8$.  The next convergent $19/12 = 1.58\overline{3}$
has $q = 12 < 3^3 = 27$.  At scale $2^3 \times 3^3$, the torus resolution is
fine enough to separate the cell error from the nearest convergent approximation,
so Baker separation~\cite{Baker1975} prevents any further splitting of cells
beyond this resolution.

\medskip\noindent\textbf{The ghost island.}
At $N = 10^{11}$, a secondary cluster of $251$ cells emerges at perpendicular
distance ${\sim}370$ from the main branch locus at $k = 729$, entirely absent
at $N = 10^{10}$ (Figures~\ref{fig:ghost} and~\ref{fig:ghost-zoom},
Table~\ref{tab:ghost}). We argue in Section~\ref{sec:discussion} that this
``ghost island'' constitutes visual proof of the Baker kick: the $352$-unit
void between the clusters is a geometric shadow of the Baker-type lower bound
on $|p - q \logb 3|$.

\begin{table}[h]
\tabcolsep=4pt
\TBL{\caption{Main cluster vs.\ ghost island at $k = 729$.}\label{tab:ghost}}
{\begin{fntable}
\begin{tabular*}{\textwidth}{@{\extracolsep{\fill}}lrr@{}}
\toprule
& \textbf{Main cluster} & \textbf{Ghost island} \\
\midrule
Total cells & $23{,}017$ & $251$ \\
Branch cells & $21{,}471$ & $161$ \\
Pure-even & $1{,}495$ & $87$ \\
Pure-odd & $51$ & $3$ \\
\midrule
Total visits & $25.1 \times 10^{12}$ & $3{,}233$ \\
Mean visits/cell & $1.09 \times 10^{9}$ & $12.9$ \\
Mean $p_{\text{odd}}$ (branch) & $0.284$ & $0.276$ \\
\botrule
\end{tabular*}
\end{fntable}}
\end{table}

\begin{figure}[t]
  \centering
  \includegraphics[width=\textwidth]{torus_heatmaps.png}
  \caption{Torus heatmaps of $p_{\text{odd}}$ at four resolutions.
  Each pixel represents a cell $(r_2, r_3)$; color encodes the
  fraction of odd steps landing in that cell. Green dots trace
  the irrational foliation $r_3 = \logb 3 \cdot r_2 \pmod{k}$;
  cyan markers indicate pure-even cells. \textbf{(a)}~$k = 81$:
  fully occupied, smooth gradient. \textbf{(b)}~$k = 144$: first
  appearance of empty and pure-parity cells. \textbf{(c)}~$k = 216$:
  transition complete, fractal-like occupancy. \textbf{(d)}~$k = 729$:
  sparse torus with ${\sim}4\%$ occupancy, branch cells concentrated
  near the irrational foliation. Data from $N = 10^{11}$.}\label{fig:heatmaps}
\end{figure}

\begin{figure}[t]
  \centering
  \includegraphics[width=\textwidth]{branch_k729_full.png}
  \caption{Full branch locus at $k = 729 = 3^6$, showing the
  main cluster (23{,}017 cells, lower band) and the ``ghost island''
  (251 cells, upper band near $r_3 \approx 440$--$461$).
  The 352-unit void between the clusters is controlled by the
  quality of the rational approximation $460/729 \approx \log_3 2$
  (error $7.2 \times 10^{-5}$). The island was entirely absent at
  $N = 10^{10}$ and emerged only in the $10^{11}$ computation.}\label{fig:ghost}
\end{figure}

\begin{figure}[t]
  \centering
  \includegraphics[width=\textwidth]{ghost_island.png}
  \caption{Zoomed view of the ghost island at $k = 729$.
  The island contains $161$ branch, $87$ pure-even, and $3$ pure-odd
  cells, with total visits $= 3{,}233$ (mean $12.9$ per cell) versus
  $1.09 \times 10^9$ in the main cluster, a visit ratio of
  $\approx 10^{-8}$. The island lies along a shifted copy of the
  irrational foliation, centered on the rational foliation
  $r_3 = (460/729)\, r_2$, and maintains
  $\bar{p}_{\text{odd}} = 0.276$. The golden mean shift constraint
  ($\mathit{oo} = 0$) is preserved in all 251~cells.}\label{fig:ghost-zoom}
\end{figure}


%% ============================================================
\section{The Spectral Gap}\label{sec:spectral}

\begin{openproblem}[Spectral gap of $\times 3$]\label{obs:spectral}
  Let $\mathcal{L}_3^{(j)}$ denote the transfer operator of the map $x \mapsto 3x$
  on the partition $\mathcal{P}_2^{(j)}$ of $\ZZ/2^j\ZZ$ into $2^{j-2}$ orbits
  (for $j \geq 3$). All non-trivial eigenvalues of $\mathcal{L}_3^{(j)}$ have
  magnitude exactly $1/3$, giving a spectral gap $\gamma = 2/3$ independent of $j$.

  The proof rests on the identity
  \[
    \prod_{k=0}^{L-1} (1 + 2\cos(3^k \theta)) = 1
  \]
  whenever $3^L \theta \equiv \theta \pmod{2\pi}$ (the product telescopes via
  $1 + 2\cos\theta = (e^{3i\theta} - 1)/(e^{i\theta} - 1)$). This forces
  $|\lambda|^L = 3^{-L}$ for each non-trivial eigenvalue, giving $|\lambda| = 1/3$.

  Verified computationally: Fourier orbit products to $j = 30$; direct matrix
  eigenvalues to $j = 10$ (both exact to $10$ significant figures).
\end{openproblem}

\begin{proposition}[Exponential mixing]\label{prop:mixing}
  For any probability distribution $p$ on $\mathcal{P}_2^{(j)}$ atoms:
  \[
    \|\mathcal{L}_3^n p - \text{uniform}\|_2 \leq (1/3)^n \cdot \|p - \text{uniform}\|_2.
  \]
  After $n$ applications of $\times 3$, the distribution on $2$-adic atoms
  is exponentially close to uniform at rate $1/3$ per step.
\end{proposition}

\begin{remark}[Multiplicative order asymmetry]
  The spectral gap has an arithmetic origin.  The order
  $\mathrm{ord}_{2^j}(3) = 2^{j-2}$ for $j \geq 3$ (index~$2$: the map
  $\times 3$ explores only \emph{half} of $(\ZZ/2^j\ZZ)^\times$), while
  $\mathrm{ord}_{3^k}(2) = \varphi(3^k) = 2 \cdot 3^{k-1}$ (index~$1$:
  $2$ is a primitive root modulo $3^k$).  Both orders have been verified
  computationally to exponent $40$.  This asymmetry means the $\times 3$ map
  mixes the $2$-adic coordinate rapidly (spectral gap $2/3$), while the
  $\times 2$ map explores the $3$-adic coordinate completely.
\end{remark}


%% ============================================================
\section{The Six Equivalences}\label{sec:equivalences}

We now state the main structural result: the Collatz conjecture is equivalent to
each of six distinct properties of the Syracuse dynamics. All equivalences are
machine-checked in Lean~4.

\begin{theorem}[Machine-checked equivalences]\label{thm:equivalences}
  The following statements are each equivalent to the Collatz conjecture
  (``every positive integer eventually reaches $1$ under iteration of $C$''):
  \begin{enumerate}
    \item[\textbf{(P1)}] \textbf{Linear $\nu_3$ bound.}
      For every $n \geq 1$, there exists $K > 0$ such that
      $3\nu_3(n,t) \leq t + K$ for all $t \geq 0$.

    \item[\textbf{(P2)}] \textbf{Finite deficit bound.}
      For every $n \geq 1$, there exists $D$ such that
      $\deficit(n,t) \leq D$ for all $t \geq 0$.

    \item[\textbf{(P3)}] \textbf{Equidistribution implies bounded deficit.}
      If the cell error sequence $\{t \logb 3 - \nu_2(n,t)\}_{t \geq 0}$
      is Weyl-equidistributed modulo $1$, then the deficit is bounded.

    \item[\textbf{(P4)}] \textbf{Cell sequence equidistribution.}
      For every odd $n \geq 1$, the sequence of $\nu_2$-residues along the
      trajectory is equidistributed in $\ZZ/k\ZZ$ for every $k$.

    \item[\textbf{(P5)}] \textbf{Syracuse valuation equidistribution.}
      For every odd $n \geq 1$, the Syracuse valuation sums are
      equidistributed on the solenoid.

    \item[\textbf{(P6)}] \textbf{Simultaneous Diophantine approximation.}
      The triple $(\log 2, \log 5, \log 7)$ satisfies a quantitative
      simultaneous approximation bound that constrains the Littlewood
      residence time of the Syracuse walk.
  \end{enumerate}
\end{theorem}

\begin{remark}
  Properties~\textbf{(P1)} and \textbf{(P2)} are elementary reformulations of
  the conjecture in terms of growth rates. Properties~\textbf{(P3)}--\textbf{(P5)}
  concern equidistribution of trajectory statistics, connecting the Collatz problem
  to ergodic theory. Property~\textbf{(P6)} is a Diophantine condition that
  bridges the number-theoretic and ergodic aspects.
\end{remark}

We now sketch the proof of each equivalence.  Full details
are machine-checked in the Lean~4 formalization.

\medskip\noindent\textbf{(P1) $\Leftrightarrow$ Collatz.}
($\Leftarrow$) If $n$ reaches $1$ at time $t^*$, then $\nu_3(n,t) = \nu_3(n,t^*)$
for $t \geq t^*$, so $3\nu_3(n,t) \leq t + 3\nu_3(n,t^*)$.
($\Rightarrow$) If $3\nu_3(n,t) \leq t + K$, then
$a_t = n \cdot 2^{-w(t)} + r(t)$ with $w(t) \geq \varepsilon t$ for
$\varepsilon = 1 - \logb 3 \cdot (K+t)/(3t) > 0$.  The correction ratio
$r(t)$ is bounded by a geometric series (Theorem~\ref{thm:r-bound} below),
so $a_t$ is eventually bounded.  By cycle elimination
(Hercher~\cite{Hercher2023} for $m \leq 91$, Baker--Steiner~\cite{Steiner1977}
for $m \geq 92$), the only cycle is $\{1,2,4\}$.
File: \texttt{Conclusion.lean}.

\medskip\noindent\textbf{(P2) $\Leftrightarrow$ Collatz.}
Immediate from (P1): $3\nu_3 \leq t + K$ is equivalent to $\deficit \leq K$
by $\deficit = 3\nu_3 - t$.  File: \texttt{DiophantineRepeller.lean}.

\medskip\noindent\textbf{(P3) $\Leftrightarrow$ Collatz.}
If the cell error sequence is Weyl-equidistributed modulo~$1$, then by
Weyl's theorem~\cite{Weyl1916} the fraction of time spent in any arc
$[a,b] \subset [0,1)$ converges to $b-a$.  Applied to the dangerous arc
(where $\val = 1$), the deficit increment averages to $2 - \mathbb{E}[\val]$,
which is negative since $\mathbb{E}[\val] \approx 1.995 > \logb 3$.
Quantitative bounds give deficit bounded.  The reverse uses
Proposition~\ref{prop:identity}: bounded trajectory implies bounded deficit
implies equidistribution on compact orbits.
File: \texttt{WeylEquidistribution.lean}.

\medskip\noindent\textbf{(P4) $\Leftrightarrow$ Collatz.}
The cell sequence $\nu_2(n,t) \bmod k$ records which column of the
$(2,3)$-torus the trajectory visits.  Equidistribution modulo every~$k$
implies that the time spent in dangerous cells has the correct density,
yielding (P3).  Conversely, bounded deficit forces the walk to stay near
the linear trajectory $w(t) \approx \varepsilon t$, and the irrationality
of $\logb 3$ (Baker) forces the cell visits to equidistribute.
File: \texttt{WeylEquidistribution.lean}.

\medskip\noindent\textbf{(P5) $\Leftrightarrow$ Collatz.}
The Syracuse valuation sum $\sum_{i=0}^{t-1} \val(3a_{s_i}+1)$ (summed over
odd steps $s_i$) records the total even-step allocation.
Equidistribution of this sum on the solenoid $\Sigma_{2,3}$ is equivalent to
equidistribution of the cell sequence (P4) by the Chinese Remainder Theorem
decomposition of cells into $2$-adic and $3$-adic components.
The proof lifts Denjoy--Koksma bounds~\cite{Herman1979} on Birkhoff sums
to the Syracuse cocycle on the skew product.
File: \texttt{SolenoidMixing.lean}.

\medskip\noindent\textbf{(P6) $\Leftrightarrow$ Collatz.}
The Littlewood residence time bounds how long a trajectory can linger in a
narrow strip of the torus.  A simultaneous approximation bound on
$(\log 2, \log 5, \log 7)$ constrains the return time to dangerous cells,
via Matveev's theorem on linear forms in three logarithms~\cite{Matveev2000}.
The reduction proceeds through the continued fraction expansion of $\logb 3$
and the residence time analysis in $2$-dimensional torus cells.
File: \texttt{LittlewoodInduction.lean}.

\begin{theorem}[Correction ratio bound]\label{thm:r-bound}
  If the walk has eventually linear drift with rate $\varepsilon > 0$
  (i.e., $w(n,t) \geq \varepsilon t$ for $t \geq T_0$), then $r(t)$ is
  eventually bounded.  Specifically, letting $\alpha > \logb 3$ denote the
  minimum safe-step valuation sum per odd step, we have $\rho = 3 \cdot 2^{-\alpha} < 1$,
  and for all $k$ sufficiently large,
  \begin{equation}\label{eq:r-bound}
    r(\tau_k + 1) \leq \frac{1}{1-\rho} + 1,
  \end{equation}
  where $\tau_k$ is the time of the $k$-th odd step.
\end{theorem}

\begin{proof}[Proof sketch]
  After the $k$-th odd step, $C(\tau_k+1) = \sum_{j=1}^{k} 3^{k-j} \cdot 2^{U_j}$
  where $U_j$ is the number of even steps before the $j$-th odd step.
  Dividing by $2^{\nu_2(\tau_k+1)}$ and using $w(t) \geq \varepsilon t$
  gives $r(\tau_k+1) \leq \sum_{j=0}^{\infty} \rho^j = 1/(1-\rho)$ plus a
  finite correction.
\end{proof}


%% ============================================================
\section{The Lean~4 Formalization}\label{sec:lean}

The formalization comprises approximately $10{,}500$ lines across $40$ files,
built on Lean~4 (v4.27.0) with Mathlib. Table~\ref{tab:axioms} lists the
$8$ axioms (published theorems used without proof) and Table~\ref{tab:sorrys}
lists the $6$ open problems (each proved equivalent to the Collatz conjecture).

\begin{table}[t]
\tabcolsep=4pt
\TBL{\caption{Axioms used in the formalization. Each corresponds to a published
theorem in transcendental number theory or ergodic theory.}\label{tab:axioms}}
{\begin{fntable}
\begin{tabular*}{\textwidth}{@{\extracolsep{\fill}}llll@{}}
\toprule
\# & Name & Reference & Used for \\
\midrule
A1 & Baker's theorem & Baker 1975 & Cycle elimination \\
A2 & Rhin's bound & Rhin 1987 & Irrationality measure \\
A3 & Hercher's bound & Hercher 2023 & $m$-cycle elimination ($m \leq 91$) \\
A4 & Weyl equidistribution & Weyl 1916 & Cell error equidistribution \\
A5 & Baker--Steiner & Steiner 1977 & Cycle length lower bound \\
A6 & Matveev three-log & Matveev 2000 & Linear forms in three logarithms \\
A7 & Denjoy--Koksma & Herman 1979 & Birkhoff sum bounds \\
A8 & Furstenberg unique erg. & Furstenberg 1961 & Skew product ergodicity \\
\botrule
\end{tabular*}
\end{fntable}}
\end{table}

\begin{table}[t]
\tabcolsep=4pt
\TBL{\caption{Open problems (sorry statements) in the formalization. Each is
machine-checked to be equivalent to the Collatz conjecture.}\label{tab:sorrys}}
{\begin{fntable}
\begin{tabular*}{\textwidth}{@{\extracolsep{\fill}}llll@{}}
\toprule
\# & Name & File & Property \\
\midrule
S1 & \texttt{nu3\_linear\_bound} & Drift.lean & (P1) \\
S2 & \texttt{finite\_deficit\_bound} & DiophantineRepeller.lean & (P2) \\
S3 & \texttt{equidist\_implies\_deficit} & WeylEquidistribution.lean & (P3) \\
S4 & \texttt{cellSeqNu2\_equidist} & WeylEquidistribution.lean & (P4) \\
S5 & \texttt{syracuseValSum\_equidist} & SolenoidMixing.lean & (P5) \\
S6 & \texttt{simult\_approx\_log2\_5\_7} & LittlewoodInduction.lean & (P6) \\
\botrule
\end{tabular*}
\end{fntable}}
\end{table}

The primary technical innovation in the formalization is the \textbf{bit-peeling
lemma} (\texttt{CarryBitScrambling.lean}), the first machine-checked proof that
the Collatz map operates as a \emph{lossy information channel}. Specifically:
$T(n) \bmod 2^k$ depends only on $n \bmod 2^{k+1}$, so each compressed step
irreversibly destroys exactly one bit of $2$-adic information. By induction, $M$
consecutive compressed steps consume $M$ bits: the output modulo $2^k$ is
determined by the input modulo $2^{k+M}$, and the erased bits are lost forever.
This gives the Collatz map a precise channel capacity: $1$ bit consumed per step,
with no possibility of reconstruction. The bit-peeling lemma is the engine behind
both the metric conflict (the erased bit shifts the cell error by $\logb 3 - 1$)
and the finite saturation of the branch locus (cells at resolution $2^a$ can only
refine $a$ times before the bit budget is exhausted).

The remaining $0$-sorry infrastructure (fully proved theorems) includes:
\begin{itemize}
  \item Cycle elimination for $1$- and $2$-cycles (\texttt{CycleElimination.lean}),
  \item Hensel attrition and the backward induction theorem (\texttt{HenselAttrition.lean}),
  \item The metric conflict: Hensel $\times$ Baker incompatibility (\texttt{MetricConflict.lean}),
  \item The fuel refill identity and divisibility transfer (\texttt{FuelDynamics.lean}),
  \item The spectral gap of $\times 3$: orbit cosine product identity (\texttt{SpectralGap.lean}),
  \item The non-coboundary theorem for the Collatz cocycle (\texttt{UniqueErgodicity.lean}),
  \item Sublinear drift from deficit boundedness (\texttt{SublinearDrift.lean}),
  \item The $E_8$ root system: $240$ roots, norm, inner product (\texttt{E8Lattice.lean}).
\end{itemize}


%% ============================================================
\section{Discussion}\label{sec:discussion}

\subsection{The information ceiling}

\begin{figure}[t]
  \centering
  \includegraphics[width=\textwidth]{convergence.png}
  \caption{Convergence of cell counts with increasing $N$.
  \textbf{(a)}~Branch cell count at $k = 216$ vs.\ checkpoint $N$:
  rapid initial growth, saturating at $21{,}632$ by $N \approx 10^8$.
  \textbf{(b)}~Multiple $k$-levels on log scale: small-$k$ levels
  saturate immediately at $k^2$; larger levels ($k = 162, 216$) take
  longer but all converge. \textbf{(c)}~Empty cell count at $k = 729$
  vs.\ $N$: slowly decreasing, but the branch cell count has already
  frozen. \textbf{(d)}~Pure-even ($1{,}582$) and pure-odd ($54$) counts
  at $k = 216$: both stabilize early, confirming robust
  classification.}\label{fig:convergence}
\end{figure}

The saturation of the branch locus at $k = 216$ suggests that the Syracuse
map admits a finite \emph{generating partition} in the sense of Adler and
Weiss~\cite{AdlerWeiss1967,AdlerWeiss1970}.  The constant complexity of
$21{,}632$ cells represents the \emph{state complexity} of the underlying
arithmetic automaton---a limit on base-reconciliation first theorized by
Cobham~\cite{Cobham1969}, whose theorem implies that sets defined
simultaneously in base~$2$ and base~$3$ can only overlap in eventually
periodic ways.  Furthermore, the spectral gap of $2/3$
(Observation~\ref{obs:spectral}) ensures that the system is \emph{uniquely
ergodic} at the partition level, a property that allows lifting from
statistical mixing of the $(2,3)$-solenoid to individual trajectories,
following the foundational work of Furstenberg~\cite{Furstenberg1961,Furstenberg1967}.
The 1967 paper provides the ``disjointness'' framework (the $2$-adic and
$3$-adic metrics do not interact), while the 1961 paper provides the unique
ergodicity that converts a spectral gap into a pointwise guarantee.

The saturation of the branch locus at $21{,}632$ cells thus suggests that the
Collatz map has \emph{finite combinatorial complexity}: beyond a certain
resolution ($k = 216$), no new qualitative behavior emerges. Every trajectory,
regardless of its starting value, visits only the same $23{,}268$ non-empty
cells on the $(2,3)$-torus. This is a strong form of ``topological finiteness''
that, to our knowledge, has not been previously observed.

\begin{table}[h]
\tabcolsep=4pt
\TBL{\caption{Computational evidence. Statistics are stable across three orders
of magnitude in starting range.}\label{tab:evidence}}
{\begin{fntable}
\begin{tabular*}{\textwidth}{@{\extracolsep{\fill}}lrrl@{}}
\toprule
\textbf{Quantity} & $N = 10^{10}$ & $N = 10^{11}$ & \textbf{Note} \\
\midrule
Total trajectory steps & $2.27 \times 10^{12}$ & $25.1 \times 10^{12}$ & \\
Mean $p_{\text{odd}}$ & $0.3324$ & $0.3325$ &
  Stable; well below $p_{\text{eq}} \approx 0.387$ \\
Branch cells ($k \geq 216$) & $21{,}632$ & $21{,}632$ & Unchanged \\
Pure-even cells & $1{,}582$ & $1{,}582$ & Unchanged \\
Pure-odd cells & $54$ & $54$ & Unchanged \\
Mean $\bar{p}_{\text{odd}}$ (branch) & $0.2843$ & $0.2843$ &
  All 4 decimal places stable \\
Mean entropy $\bar{H}$ & $0.8380$ & $0.8380$ & All 4 decimal places stable \\
``$11$'' SFT violations & $0$ & $0$ &
  Over $25.1 \times 10^{12}$ transitions \\
\botrule
\end{tabular*}
\end{fntable}}
\end{table}

We emphasize that saturation of the branch locus does \emph{not} imply convergence.
A finite graph can support non-terminating paths, and the transition structure
within the $21{,}632$ branch cells is not Markovian (transitions depend on the
full value of $n$, not just its cell). The gap between ``finite static complexity''
and ``universal convergence'' is precisely the content of properties
(P1)--(P6).

\begin{figure}[t]
  \centering
  \includegraphics[width=\textwidth]{transitions.png}
  \caption{Transition and shadow analysis at $k = 729$, $N = 10^{11}$.
  \textbf{Top left:} Transition valence map, each non-empty cell is colored
  by the number of distinct transition types (ee, eo, oe) it exhibits; the
  ghost island appears as a secondary band near $r_3 \approx 440$--$460$.
  \textbf{Top center:} ``$11$'' successor analysis verifying the golden mean
  shift at each cell. \textbf{Top right:} Shadow offset spatial map, red
  cells are closer to the rational foliation $460/729$, blue cells closer
  to the irrational foliation $\logb 3$. \textbf{Bottom right:} Shadow offset
  distribution with gap near zero and asymmetric shoulder at $+0.007$ from the
  ghost island.}\label{fig:transitions}
\end{figure}

\subsection{The ghost island: visual proof of the Baker kick}

At $N = 10^{11}$, a secondary cluster of $251$ cells emerges at perpendicular
distance ${\sim}370$ from the main branch locus at $k = 729$, entirely absent
at $N = 10^{10}$ (Figures~\ref{fig:ghost} and~\ref{fig:ghost-zoom}). This
\emph{ghost island} is a direct, visible consequence of Baker's theorem in action.

The island lies along the rational approximant $460/729 \approx \log_3 2$
(error $7.2 \times 10^{-5}$), and the $352$-unit void separating it from the
main cluster is a geometric shadow of the Baker-type lower bound
$|p - q \logb 3| > C/q^{\kappa}$. No cells can exist in the void: Baker
separation forbids it. The ghost island is populated only by trajectories that
reach the second-nearest rational approximant of $\logb 3$, an event so rare
that it requires $> 10^{10}$ starting values to observe even a single visit.

The ghost island is significant because it provides \emph{visual proof} that
the ``Baker kick'' is real and operates at exactly the scale predicted by the
theory. Each danger step shifts the cell error by $\logb 3 - 1 \approx 0.585$,
and this shift accumulates until the trajectory is ejected from the main cluster
entirely. The $352$-unit void is the minimum ejection distance, set by the
quality of the best rational approximation at scale $k = 729$. The ghost island's
statistics ($\bar{p}_{\text{odd}} = 0.276$, identical SFT constraint, cell type
distribution proportional to the main cluster) confirm that it is governed by the
same dynamics, merely displaced by one Baker kick.

\subsection{The gap}

The honest assessment: each of the six open problems (P1)--(P6) is equivalent
to the Collatz conjecture. No reformulation, however elegant, changes this.
The contribution of this paper is not a proof of the conjecture, but a
\emph{complete structural map} of its irreducible content.

The structure decomposes into ``microscopic'' and ``macroscopic'' layers:

\begin{table}[h]
\tabcolsep=4pt
\TBL{\caption{Three structural forces constraining Syracuse trajectories.}\label{tab:forces}}
{\begin{fntable}
\begin{tabular*}{\textwidth}{@{\extracolsep{\fill}}llp{0.42\textwidth}l@{}}
\toprule
\textbf{Force} & \textbf{Metric} & \textbf{Effect} & \textbf{Status} \\
\midrule
Hensel attrition & $2$-adic &
  Danger runs of length $d$ require $2^{d+1} \mid (n+1)$;
  density $2^{-d}$ among odd numbers & Proved \\
Baker separation & Archimedean &
  $|\varepsilon(a,b)| > C/\max^{\kappa}$;
  cells cannot cluster near critical line & Proved (mod A1) \\
Weyl equidistribution & Ergodic &
  Cell visits uniform mod~$k$;
  no systematic avoidance of safe cells & \textbf{Open} ($\equiv$ Collatz) \\
\botrule
\end{tabular*}
\end{fntable}}
\end{table}

\noindent The first two forces (``microscopic'') are fully proved. Hensel attrition shows
that individual danger runs are exponentially rare, and Baker separation shows
that they are spatially isolated on the torus. Carry-bit scrambling gives
exact $1$-bit erosion per compressed step, so the Collatz map genuinely
``forgets'' information.

The third force (``macroscopic'') is the open problem: proving that the
microscopic mixing, which is entirely deterministic, not random, implies
trajectory-level equidistribution of cell visits. This is the irreducible
content of the Collatz conjecture, common to all six formulations (P1)--(P6).

\subsection{Connection to Tao's approach}

\begin{table}[t]
\tabcolsep=4pt
\TBL{\caption{Comparison of approaches to the Collatz conjecture.}\label{tab:comparison}}
{\begin{fntable}
\begin{tabular*}{\textwidth}{@{\extracolsep{\fill}}lll@{}}
\toprule
& \textbf{Tao~\cite{Tao2022}} & \textbf{This paper} \\
\midrule
\emph{Result type} & Density (almost all orbits) & Equivalence (all orbits) \\
\emph{Strength} & Almost bounded values & Full conjecture $\Leftrightarrow$ 6 properties \\
\emph{Map} & Syracuse $\Syr$ on $\ZZ/3^n\ZZ$ & Syracuse on $(2,3)$-solenoid \\
\emph{Key tool} & Characteristic functions & Baker separation + Hensel attrition \\
\emph{Equidistribution} & Syracuse r.v.'s on $\ZZ/3^n\ZZ$ & Cell sequences on torus \\
\emph{Mixing source} & Entropy of Syracuse r.v.'s & Spectral gap $2/3$ of $\times 3$ \\
\emph{Gap} & Logarithmic density $\to$ natural & Partition mixing $\to$ trajectory \\
\emph{Computation} & None & $2.5 \times 10^{13}$ steps, GPU \\
\emph{Formalization} & None & $10{,}500$ lines Lean~4 \\
\emph{Cycle elimination} & Not needed (density result) & Baker + Hercher ($m \leq 91$) \\
\botrule
\end{tabular*}
\end{fntable}}
\end{table}

Tao's method~\cite{Tao2022} controls the \emph{typical} behavior of Syracuse
random variables via characteristic function bounds on $\ZZ/3^n\ZZ$. His result, that
almost all orbits (in the sense of logarithmic density) attain almost bounded
values, is the deepest partial result toward the conjecture. Our framework
complements this in three ways:
\begin{enumerate}
  \item \textbf{Individual vs.\ typical.} Tao's approach proves that most starting
    values behave well; ours identifies what would be needed for \emph{every}
    starting value. The six equivalences (P1)--(P6) constitute a precise
    ``wishlist'' for an individual-trajectory proof.
  \item \textbf{Mixing mechanism.} Tao uses the entropy of Syracuse random variables
    on $\ZZ/3^n\ZZ$; we identify the spectral gap $\gamma = 2/3$ of the $\times 3$
    transfer operator on $2$-adic partitions (Observation~\ref{obs:spectral}) as
    the structural source of mixing. Both approaches ultimately rely on the
    multiplicative independence of $2$ and $3$.
  \item \textbf{Machine-checked reduction.} The equivalence chain is verified in
    Lean~4, providing a formal ``map'' of the irreducible content. Any future
    proof of the conjecture must (implicitly or explicitly) establish one of
    (P1)--(P6).
\end{enumerate}

The spectral gap $\gamma = 2/3$ for $\times 3$ on $2$-adic partitions
provides the mixing rate at the partition level.
Lifting this to trajectory-level equidistribution remains the central challenge:
the gap between Tao's ``almost all'' and our ``every trajectory'' is precisely
the gap between statistical and pointwise ergodic theorems for the Syracuse
dynamics.


\begin{Backmatter}

\paragraph{Acknowledgments}
The Lean~4 formalization uses the Mathlib library. Computations were performed
on an Intel Core Ultra~9 275HX (24 cores) with 128\,GB DDR5 RAM and an
NVIDIA RTX~5090 (24\,GB GDDR7) for GPU-accelerated branch locus computation.

\paragraph{Competing interests}
None.

\paragraph{Data availability statement}
The Lean~4 formalization, C source code for all computations, and computed
branch locus data are available at
\url{https://github.com/johnjanik/syracuse-confinement}.

\begin{thebibliography}{99}

\bibitem{AdlerWeiss1967}
R.~L.~Adler and B.~Weiss, Entropy, a complete metric invariant for
automorphisms of the torus, \emph{Proc.\ Nat.\ Acad.\ Sci.\ U.S.A.}
\textbf{57} (1967), 1573--1576.

\bibitem{AdlerWeiss1970}
R.~L.~Adler and B.~Weiss, \emph{Similarity of Automorphisms of the Torus},
Memoirs of the American Mathematical Society, No.~98, AMS, 1970.

\bibitem{Baker1975}
A.~Baker, \emph{Transcendental Number Theory}, Cambridge University Press, 1975.

\bibitem{Barina2025}
D.~Barina, Improved verification limit for the convergence of the Collatz
conjecture, \emph{J.\ Supercomput.} \textbf{81} (2025), 810.

\bibitem{Cobham1969}
A.~Cobham, On the base-dependence of sets of numbers recognizable by finite
automata, \emph{Math.\ Systems Theory} \textbf{3} (1969), 186--192.

\bibitem{Furstenberg1961}
H.~Furstenberg, Strict ergodicity and transformation of the torus,
\emph{Amer.\ J.\ Math.} \textbf{83} (1961), 573--601.

\bibitem{Furstenberg1967}
H.~Furstenberg, Disjointness in ergodic theory, minimal sets, and a problem
in Diophantine approximation, \emph{Math.\ Systems Theory} \textbf{1} (1967),
1--49.

\bibitem{HeuleAaronsonYolcu2023}
M.~Heule, S.~Aaronson, and E.~Yolcu, An automated approach to the Collatz
conjecture, \emph{J.\ Automat.\ Reason.} (2023).

\bibitem{Hercher2023}
J.~Hercher, There are no Collatz $m$-cycles with $m \leq 91$,
\emph{J.\ Integer Seq.} \textbf{26} (2023).

\bibitem{Herman1979}
M.~Herman, Sur la conjugaison diff\'erentiable des diff\'eomorphismes du cercle
\`a des rotations, \emph{Publ.\ Math.\ IH\'ES} \textbf{49} (1979), 5--233.

\bibitem{Lagarias2003}
J.~C.~Lagarias, The $3x+1$ problem: An annotated bibliography (1963--1999),
\texttt{arXiv:math/0309224}, 2003.

\bibitem{Lagarias2010}
J.~C.~Lagarias (ed.), \emph{The Ultimate Challenge: The $3x+1$ Problem},
AMS, 2010.

\bibitem{Matveev2000}
E.~M.~Matveev, An explicit lower bound for a homogeneous rational linear form
in the logarithms of algebraic numbers. II, \emph{Izv.\ Ross.\ Akad.\ Nauk
Ser.\ Mat.} \textbf{64} (2000), 125--180.

\bibitem{Rhin1987}
G.~Rhin, Approximants de Pad\'e et mesures effectives d'irrationalit\'e,
\emph{S\'eminaire de Th\'eorie des Nombres, Paris 1985--86}, Progress in Math.
\textbf{71} (1987), 155--164.

\bibitem{ShakibaeiAsli2026}
B.~Shakibaei Asli, An explicit near-conjugacy between the Collatz map and a
circle rotation, \texttt{arXiv:2601.04289}, 2026.

\bibitem{SimonsDeWeger2005}
J.~Simons and B.~M.~M.~de Weger, Theoretical and computational bounds for
$m$-cycles of the $3n+1$ problem, \emph{Acta Arith.} \textbf{117} (2005),
51--70.

\bibitem{Steiner1977}
R.~P.~Steiner, A theorem on the Syracuse problem, in \emph{Proceedings of the
7th Manitoba Conference on Numerical Mathematics}, 1977, pp.~553--559.

\bibitem{Tao2022}
T.~Tao, Almost all orbits of the Collatz map attain almost bounded values,
\emph{Forum Math.\ Pi} \textbf{10} (2022), e12.

\bibitem{Weyl1916}
H.~Weyl, \"Uber die Gleichverteilung von Zahlen mod.\ Eins,
\emph{Math.\ Ann.} \textbf{77} (1916), 313--352.

\end{thebibliography}

\end{Backmatter}

\end{document}
